\subsection{微分と一次近似}
微分可能性は,
\[
 f(a+h)=f(a)+f'(a)h+h\varepsilon(h),
 \qquad \varepsilon(h)\to0\quad(h\to0)
\]
と書き直せる.実際,\,$h\ne0$ に対して $\varepsilon(h)=\frac{f(a+h)-f(a)}h-f'(a)$ と置けば,微分係数の定義から $\varepsilon(h)\to0$ となる.\,$\varepsilon(h)$ の値は $h$ によって変わり,正にも負にもなり得る.最後の $h\varepsilon(h)$ は,\,$h$ で割っても $0$ に近づく誤差である.右辺の初めの2項 $f(a)+f'(a)h$ は,点 $(a,f(a))$ における接線の $x=a+h$ での値にあたる.つまり,\,$h$ が小さいとき,関数の値と接線の値の差は $h$ よりもさらに速く $0$ に近づくのである.この形は合成関数の微分を証明するときにも役立つ.
